% TODO checklist:
% - [x] Inspected src/jobs/models.py (SSEEvent model)
% - [x] Inspected src/jobs/interfaces.py (EventStore interface)
% - [x] Inspected db/redis_cache/job_store.py (RedisEventStore implementation)
% - [x] Inspected src/jobs/metrics.py (SSE metrics)
% - [x] Inspected src/routers/jobs.py (SSE endpoint pattern)
% - [x] Analyzed Q&A.md Q38 for event model details

\section{Job Events and SSE Streaming}
\label{sec:job-events}

Real-time progress updates are essential for user experience when executing long-running jobs. The Cineca Agentic Platform implements Server-Sent Events (SSE) streaming with a ring buffer architecture that supports client reconnection and event replay. This section describes the event model, storage mechanism, and streaming protocol implementation.

\subsection{Design Rationale}

The SSE streaming subsystem addresses several requirements:

\begin{itemize}
    \item \textbf{Real-Time Updates}: Clients receive immediate notification of status transitions without polling.
    
    \item \textbf{Resume Capability}: Disconnected clients can resume from their last received event using the \texttt{Last-Event-ID} header.
    
    \item \textbf{Bounded Storage}: Event history is limited to prevent unbounded memory growth while retaining recent history for reconnection.
    
    \item \textbf{Protocol Compliance}: Events follow the W3C Server-Sent Events specification for broad client compatibility.
\end{itemize}

\subsection{SSE Event Model}

The \texttt{SSEEvent} model in \texttt{src/jobs/models.py} represents a single event in the job event stream:

\begin{lstlisting}[language=Python, caption={SSEEvent model definition.}]
class SSEEvent(BaseModel):
    """Server-Sent Event for job status updates."""

    event_id: int = Field(..., description="Monotonic event sequence number")
    event_type: str = Field(..., description="Event type: 'status', 'end', 'error'")
    data: dict[str, Any] = Field(..., description="Event payload")
    timestamp: datetime = Field(
        default_factory=datetime.utcnow,
        description="Event emission time (UTC)"
    )

    def to_sse_format(self) -> str:
        """Format as SSE wire protocol."""
        lines = [
            f"id: {self.event_id}",
            f"event: {self.event_type}",
            f"data: {json.dumps(self.data)}",
            "",  # SSE requires blank line terminator
        ]
        return "\n".join(lines)
\end{lstlisting}

\subsection{Event Types}

The job system emits several event types throughout the job lifecycle:

\begin{table}[ht]
\centering
\caption{SSE event types and their semantics.}
\label{tab:sse-event-types}
\small
\begin{tabular}{llp{6cm}}
\hline
\textbf{Event Type} & \textbf{Trigger} & \textbf{Data Payload} \\
\hline
\texttt{status} & Status transition & \texttt{\{from, to, timestamp\}} \\
\texttt{progress} & Progress update & \texttt{\{percent, step, message\}} \\
\texttt{heartbeat} & Keep-alive & \texttt{\{timestamp\}} \\
\texttt{end} & Terminal state reached & \texttt{\{status, result?, error?\}} \\
\texttt{error} & Error during streaming & \texttt{\{message, code\}} \\
\hline
\end{tabular}
\end{table}

\subsection{Ring Buffer Architecture}

Events are stored in a ring buffer implemented using Redis LIST data structures. The ring buffer provides bounded storage with automatic eviction of oldest events:

\begin{figure}[ht]
\centering
\begin{tikzpicture}[
    event/.style={rectangle, draw, minimum width=1.5cm, minimum height=0.8cm, font=\small},
    arrow/.style={->, >=stealth, thick}
]
    % Ring buffer visualization
    \node[event, fill=gray!30] (e1) at (0,0) {Event 1};
    \node[event, fill=gray!30] (e2) at (2,0) {Event 2};
    \node[event, fill=blue!20] (e3) at (4,0) {Event 3};
    \node[event, fill=blue!20] (e4) at (6,0) {Event 4};
    \node[event, fill=green!20] (e5) at (8,0) {Event 5};
    
    \node[above=0.3cm of e1, font=\scriptsize] {oldest};
    \node[above=0.3cm of e5, font=\scriptsize] {newest};
    
    \draw[arrow] (9.5,0) -- node[right, font=\scriptsize] {LPUSH} (8.8,0);
    \draw[arrow] (-0.8,0) -- node[left, font=\scriptsize] {LTRIM evicts} (-1.5,0);
    
    % Legend
    \node[below=1cm of e3, font=\small] {Ring size = 5 (configurable via SSE\_RING\_SIZE)};
\end{tikzpicture}
\caption{Ring buffer event storage with LPUSH/LTRIM operations.}
\label{fig:ring-buffer}
\end{figure}

The Redis implementation uses atomic LPUSH + LTRIM to maintain the ring buffer:

\begin{lstlisting}[language=Python, caption={Ring buffer append with atomic trim.}]
async def append(self, job_id: str, event: SSEEvent, ring_size: int) -> None:
    """Append event to ring buffer, capping at ring_size."""
    redis = await get_async_redis()
    events_key = f"job:{job_id}:events"
    event_json = event.to_storage_json()

    async with redis.pipeline(transaction=True) as pipe:
        # Prepend event (newest first)
        pipe.lpush(events_key, event_json)
        # Trim to ring_size (keeps indices 0 to ring_size-1)
        pipe.ltrim(events_key, 0, ring_size - 1)
        # Refresh TTL to match job expiry
        pipe.expire(events_key, self._ttl_seconds)
        await pipe.execute()
\end{lstlisting}

\subsection{Monotonic Event IDs}

Each event receives a monotonically increasing ID within the job scope, enabling clients to detect gaps and request replay. The ID is generated using Redis INCR for atomic increment:

\begin{lstlisting}[language=Python, caption={Monotonic event ID generation.}]
async def get_next_event_id(self, job_id: str) -> int:
    """Get next event ID (atomic increment)."""
    redis = await get_async_redis()
    seq_key = f"job:{job_id}:event_seq"

    # INCR is atomic and returns the new value
    event_id = await redis.incr(seq_key)

    # Set TTL if this is the first event
    if event_id == 1:
        await redis.expire(seq_key, self._ttl_seconds)

    return event_id
\end{lstlisting}

\subsection{Event Replay for Reconnection}

When a client reconnects with a \texttt{Last-Event-ID} header, the server replays missed events from the ring buffer:

\begin{lstlisting}[language=Python, caption={Event replay from Last-Event-ID.}]
async def replay_from(self, job_id: str, last_event_id: int) -> list[SSEEvent]:
    """Retrieve events with event_id > last_event_id."""
    redis = await get_async_redis()
    events_key = f"job:{job_id}:events"

    # Get all events from LIST
    event_jsons = await redis.lrange(events_key, 0, -1)

    # Parse and filter to those after last_event_id
    events = []
    for event_json_bytes in reversed(event_jsons):  # Chronological order
        event_json = event_json_bytes.decode("utf-8")
        event_dict = json.loads(event_json)

        event = SSEEvent(
            event_id=event_dict["event_id"],
            event_type=event_dict["event_type"],
            data=event_dict["data"],
        )

        if event.event_id > last_event_id:
            events.append(event)

    return sorted(events, key=lambda e: e.event_id)
\end{lstlisting}

If the ring buffer has rotated past the requested \texttt{Last-Event-ID}, the server emits a comment indicating the gap:

\begin{lstlisting}[language=Python, caption={Gap detection in SSE streaming.}]
# If requested last_event_id is not in buffer, emit gap comment
if last_event_id > 0 and not replayed_events:
    yield ": gap detected, some events may have been lost\n\n"
    record_sse_gap(backend="redis")
\end{lstlisting}

\subsection{SSE Endpoint Implementation}

The SSE streaming endpoint in \texttt{src/routers/jobs.py} implements the full streaming protocol:

\begin{lstlisting}[language=Python, caption={SSE streaming endpoint skeleton.}]
@router.get("/{job_id}/events")
async def job_events(
    job_id: str,
    request: Request,
    last_event_id: int | None = Header(None, alias="Last-Event-ID"),
):
    """Stream job events via Server-Sent Events."""
    
    async def event_generator():
        # Track active connection
        with track_sse_connection(backend="redis"):
            # Replay missed events if reconnecting
            if last_event_id:
                replayed = await event_store.replay_from(job_id, last_event_id)
                for event in replayed:
                    yield event.to_sse_format()
                if replayed:
                    record_sse_resume(backend="redis")

            # Poll for new events until job terminates
            while True:
                job = await job_store.get(job_id)
                if not job:
                    yield 'event: end\ndata: {"status":"disappeared"}\n\n'
                    break

                if job.status.is_terminal:
                    yield f'event: end\ndata: {{"status":"{job.status.value}"}}\n\n'
                    record_sse_terminal(backend="redis", status=job.status.value)
                    break

                # Emit heartbeat to keep connection alive
                yield f': heartbeat {datetime.utcnow().isoformat()}\n\n'
                record_sse_heartbeat(backend="redis")
                
                await asyncio.sleep(1.0)  # Poll interval

    return StreamingResponse(
        event_generator(),
        media_type="text/event-stream",
        headers={
            "Cache-Control": "no-cache",
            "Connection": "keep-alive",
            "X-Accel-Buffering": "no",  # Disable nginx buffering
        },
    )
\end{lstlisting}

\subsection{SSE Wire Protocol}

The SSE wire protocol follows the W3C specification with the following format:

\begin{verbatim}
id: 42
event: status
data: {"job_id": "abc-123", "from": "queued", "to": "running"}

id: 43
event: progress
data: {"percent": 50, "step": 5, "message": "Processing..."}

id: 44
event: end
data: {"status": "finished", "result": {...}}
\end{verbatim}

Each event block is terminated by a double newline (\texttt{\textbackslash n\textbackslash n}), and comments (lines starting with \texttt{:}) are used for heartbeats and gap notifications.

\subsection{SSE Metrics}

Comprehensive metrics track SSE connection health and usage patterns:

\begin{lstlisting}[language=Python, caption={SSE-specific Prometheus metrics.}]
# Active connection tracking
sse_connections_active = Gauge(
    "sse_connections_active", "Active SSE connections", ["backend"]
)

# Resume functionality
sse_resume_hits_total = Counter(
    "sse_resume_hits_total", "Successful Last-Event-ID resumes", ["backend"]
)

# Gap detection
sse_gap_events_total = Counter(
    "sse_gap_events_total", "SSE ring buffer gaps (no backlog)", ["backend"]
)

# Keep-alive tracking
sse_heartbeat_total = Counter(
    "sse_heartbeat_total", "SSE heartbeats sent", ["backend"]
)

# Terminal event tracking
sse_terminal_events_total = Counter(
    "sse_terminal_events_total", "SSE terminal end events",
    ["backend", "status"]  # finished, failed, cancelled, disappeared
)
\end{lstlisting}

The \texttt{track\_sse\_connection} context manager maintains the active connection gauge:

\begin{lstlisting}[language=Python, caption={SSE connection tracking context manager.}]
def track_sse_connection(backend: str):
    """Context manager to track SSE connection lifecycle."""
    class SSEConnectionTracker:
        def __enter__(self):
            sse_connections_active.labels(backend=backend).inc()
            return self

        def __exit__(self, exc_type, exc_val, exc_tb):
            sse_connections_active.labels(backend=backend).dec()

    return SSEConnectionTracker()
\end{lstlisting}

\subsection{Client Integration}

Clients can connect to the SSE stream using standard EventSource APIs or HTTP clients with streaming support:

\begin{lstlisting}[language=JavaScript, caption={JavaScript EventSource client example.}]
const eventSource = new EventSource('/v1/jobs/abc-123/events', {
    headers: { 'Authorization': 'Bearer <token>' }
});

eventSource.addEventListener('status', (event) => {
    const data = JSON.parse(event.data);
    console.log(`Job transitioned: ${data.from} -> ${data.to}`);
});

eventSource.addEventListener('end', (event) => {
    const data = JSON.parse(event.data);
    console.log(`Job finished with status: ${data.status}`);
    eventSource.close();
});

eventSource.onerror = (error) => {
    // Browser automatically reconnects with Last-Event-ID
    console.log('Connection error, reconnecting...');
};
\end{lstlisting}

\subsection{Configuration Parameters}

Table~\ref{tab:sse-config} summarizes SSE-related configuration parameters.

\begin{table}[ht]
\centering
\caption{SSE streaming configuration parameters.}
\label{tab:sse-config}
\small
\begin{tabular}{lll}
\hline
\textbf{Parameter} & \textbf{Default} & \textbf{Description} \\
\hline
\texttt{SSE\_RING\_SIZE} & 100 & Maximum events retained per job \\
\texttt{SSE\_POLL\_INTERVAL} & 1.0 & Polling interval for new events (seconds) \\
\texttt{SSE\_HEARTBEAT\_INTERVAL} & 15.0 & Keep-alive heartbeat interval (seconds) \\
\hline
\end{tabular}
\end{table}

